\section{双线谱系：两条脉络的编年并置}\label{sec:genealogy}

\subsection{并置的理由}

把两条文献脉络放在同一张时间轴上，是检验引言论断的第一个动作。若“两策略共享同一源头”成立，那么在早期（1920--1980）的文献中，我们应当看到同一批作者用同一套统计工具同时讨论两个方向，而不是两条互不相干的溪流。事实正是如此。

\subsection{共同的数学源头（1900--1930）}

现代金融时间序列的数学起点通常追溯到 Bachelier(1900) 的博士论文，他把价格变动建模为布朗运动\cite{bachelier1900}。这一建模选择本身就规定了后续争论的形态：一旦价格服从带独立增量的随机游走，收益就既无自相关、也无均值回复，任何基于历史价格的策略在理论上都无法获利。此后近一个世纪的实证研究，本质上是不断检验这个零假设在多大程度上被违反。

Uhlenbeck 与 Ornstein(1930) 提出的过程给出了“违反”的最简便形式：一个带线性回复力的扩散过程，偏离均值越远，被拉回的倾向越强\cite{uo1930}。这个后来被称为 OU 过程的对象，在数学上是均值回归策略的母体——它同时规定了回归强度（半衰期）与平稳分布，把“价格会回来”这句话变成了可估参数。它与随机游走的差别只有一个漂移项：$dX_t=\theta(\mu-X_t)dt+\sigma dW_t$ 中的 $\theta$ 是正、是零、还是负。\textbf{这一个符号，就是本文全篇要反复回到的地方。}

\subsection{效率市场的检验期（1930--1970）：同一位作者，两个方向}

1930 至 1970 年是两条脉络共用工具的时期。

Cowles(1933) 与 Cowles \& Jones(1937) 检验专业预测者和技术交易规则能否跑赢市场，方法是对收益序列做自相关与游程检验，结论偏负面\cite{cowles1933,cowlesjones1937}。这项工作的重要性不在于结论，而在于它确立了“用序列相关检验可预测性”的方法论范式——而这正是后来动量与反转研究共同使用的工具。同一路径上，Fama 与 Blume(1966) 对“过滤规则”（filter rules）的检验给出了更直接的交易含义：在扣除交易成本后，机械的顺势规则无法战胜买入持有\cite{famablume1966}。这一结论在三十年后被 Brock 等(1992) 的移动平均与突破规则检验所动摇\cite{brock1992}，而后者又被 Sullivan 等(1999) 用数据窥探校正部分推翻\cite{sullivan1999}——这条“发现—质疑—再质疑”的链条是趋势研究中最典型的方法论演进样本。

Kendall(1953) 对英国商品与股票价格的时间序列分析，进一步支持了价格近似随机游走的判断\cite{kendall1953}。Fama(1970) 的效率市场综述把上述结果收拢为一个可检验的框架：弱式有效意味着历史价格对预测未来收益无增量信息\cite{fama1970}。Levy(1967) 的相对强弱选股准则则提供了一个早期的经验反例：按过去涨幅排序的组合在随后一段时间表现更优\cite{levy1967}——这是被后来的截面动量文献反复引用的先声。

这段历史的微妙之处在于：\emph{它同时是动量与反转的先验否定}。如果市场是弱式有效的，“追涨”与“买跌”都无效；两者的合法性必须等到效率市场的经验反例足够多之后才被重新讨论。

\subsection{分岔点：1985--1993}

真正的分岔发生在 1980 年代中期到 1990 年代初，而且两条脉络几乎同时破土。

\textbf{反转一侧}。Shiller(1981) 与 Summers(1986) 从价格波动过大这一现象出发，论证价格可以长期偏离基本面\cite{shiller1981,summers1986}。De Bondt \& Thaler(1985) 给出了最直接的交易含义：以 1926--1982 年数据构造形成期 36 个月的输家与赢家组合，发现输家随后跑赢赢家，被解释为过度反应后的回归\cite{debondt1985}。Fama \& French(1988) 用回归分离股价的永久与暂时成分，发现长周期负自相关，为“均值回归在长尺度上存在”提供了更严格的计量证据\cite{famafrench1988}。Poterba \& Summers(1988) 以方差比检验在多个国家与更长的历史样本上得到类似结论\cite{poterbsummers1988}。这一路径随后在国际样本上被反复检验：Richards(1997) 检验各国股指的赢家——输家反转能否用风险解释\cite{richards1997}，Balvers 等(2000) 把参数化反向策略推广到 18 个国家指数并给出可交易的截面对冲方案\cite{balvers2000}。

\textbf{动量一侧}。几乎同时，Jegadeesh \& Titman(1993) 报告了方向相反的结果：以 3--12 个月为形成期、3--12 个月为持有期，买入赢家卖出输家的组合在随后一年获得显著超额收益\cite{jegadeesh1993}。这与长周期反转并不矛盾——两者作用在不同的时间尺度上。Lehmann(1990) 与 Jegadeesh(1990) 则报告了更短尺度（周频、月频）上的负自相关，即短期反转\cite{lehmann1990,jegadeesh1990}。Chan 等(1996) 进一步检验动量的收益来源，指出盈利惯性与信息扩散无法完全解释，价格动量本身独立有效\cite{chan1996}。

至此，经验事实已经显露出它的真实形状：\textbf{收益自相关随尺度变化符号}。短尺度为负、中尺度为正、长尺度为负。趋势与反转不是两个事实，而是同一个自相关函数在不同横坐标上的取值。

\subsection{两条脉络的独立生长（1990--2010）}

意识到“尺度决定归属”之后，两条脉络并未合流，反而各自向纵深发展，形成两套几乎不交叉的技术语言。

\textbf{反转一侧走向统计套利}。Lo \& MacKinlay(1990) 澄清了反转利润的来源问题：许多被报告的反转收益来自股票之间的交叉自相关（大市值滞后带动小市值），而非个股自身的过度反应\cite{lomackinlay1990}。这一批评把反转研究推向更精细的技术路线。Engle \& Granger(1987) 的协整框架为此提供了计量基础\cite{englegranger1987}：只有当价差平稳（协整）时，“价差会回归”才是可检验的命题。此后配对交易成为主力范式——Gatev 等(2006) 用 1962--2002 年 CRSP 日数据，以最小化标准化价差平方和的距离法选对，报告年化超额收益约 11\%\cite{gatev2006}；Avellaneda \& Lee(2010) 用因子与主成分剥离残差构建统计套利组合\cite{avellaneda2010}；Elliott 等(2005) 把 Kalman 滤波引入配对价差的在线估计\cite{elliott2005}。Do \& Faff(2012) 随后发现配对交易利润随时间显著衰减，早期高收益不能外推\cite{dofaff2012}。

\textbf{动量一侧走向跨资产因子}。Rouwenhorst(1998) 在 12 个欧洲市场确认动量的普遍性\cite{rouwenhorst1998}；Moskowitz \& Grinblatt(1999) 发现行业动量解释了个股动量的大部分\cite{moskowitzgrinblatt1999}；Jegadeesh \& Titman(2001) 用 1990 年代的数据确认动量持续存在，并倾向于行为解释\cite{jegadeesh2001}。Fung \& Hsieh(1997,2001) 把 CTA 的收益刻画为“回溯跨式”式的凸性暴露，使趋势跟踪第一次获得了可因子化的表述\cite{funghsieh1997,funghsieh2001}。Moskowitz、Ooi \& Pedersen(2012) 提出的时间序列动量把趋势跟踪彻底学术化：在 58 个期货合约上，过去 12 个月超额收益正向预测未来收益，且几乎每个合约都盈利\cite{mop2012}。Asness 等(2013) 进一步证明价值与动量在八类资产上普遍存在且负相关\cite{asness2013}。

两套语言的差异此时已经固化：反转研究的关键词是\emph{协整、平稳性、残差、阈值}；动量研究的关键词是\emph{形成期、持有期、波动率缩放、横截面排序}。这一点将直接导致后文讨论的“入场门槛”差异。

跨资产层面，这条脉络还向外汇与商品延伸：Menkhoff 等(2012) 系统检验外汇动量策略，发现收益高但换手与交易成本显著\cite{menkhoff2012}；Erb \& Harvey(2006) 在商品期货上检验展期收益与动量择时的作用\cite{erbharvey2006}。技术分析的学术化则以 Lo 等(2000) 为代表，该文用计算算法与统计推断为形态识别提供了可检验的实现\cite{lo2000}。Baltas 与 Kosowski(2020) 后来把时间序列动量的收益拆解为波动率估计、交易规则与两两相关性三部分，是该领域最系统的技术性梳理\cite{baltaskosowski2020}。

\subsection{共同的批判期（2015--2020）}

数据窥探的批评其实早有预备：Sullivan 等(1999) 对 7846 条技术交易规则施加 bootstrap 校正，发现样本外超额收益不再显著\cite{sullivan1999}；White(2000) 提出 Reality Check\cite{white2000}。但真正的转折发生在 2015 年之后：Harvey 等(2016) 统计 316 个已发表因子，论证常规 $t>2$ 门槛过低，新因子应达到 $t>3$ 左右\cite{hlz2016}；McLean \& Pontiff(2016) 量化了发表后衰减——组合收益在样本外下降 26\%，发表后累计下降约 58\%\cite{mcleanpontiff2016}。

在这一审查下，两条脉络各自被点名。动量一侧，Huang 等(2020) 逐资产回归后质疑时间序列动量的稳健性，指出其表现与“按历史样本均值投资”的简单策略几乎无差别\cite{huang2020}；Daniel \& Moskowitz(2016) 则揭示了动量的深度左尾——在恐慌状态下的动量崩溃\cite{danielmoskowitz2016}。反转一侧，Do \& Faff(2012) 的利润衰减证据\cite{dofaff2012}、以及 2007 年 8 月量化危机的经验研究\cite{khandanilo2011}，指向反转策略在拥挤与去杠杆环境下的集体失败。

\subsection{百年尺度的证据与当代回归（2017--至今）}

Hurst、Ooi \& Pedersen(2017) 手工收集 1880 年以来的数据，把趋势跟踪的回溯拉到 138 年，发现每个十年均录得正收益，且在 60/40 组合的十次最大回撤中有八次为正\cite{hop2017}。这是“趋势=稳健因子”叙事最常被引用的证据，也构成对“样本期依赖”批评的直接回应。

与此同时，两策略在方法上的重新接近开始出现。Wood 等(2021) 的“慢动量加快反转”用深度学习与变点检测把两个方向融合进同一个交易规则\cite{wood2021}。这一尝试的动机恰恰是本文的核心论断：既然二者共享信号，那么把它们当作\emph{同一信号在不同尺度上的两个分量}来建模，在工程上是自然的。

\subsection{编年并置说明了什么}

把两条脉络放在一起看，有三个事实浮现。

\textbf{其一，源头是共同的}。OU 过程同时孕育了两者；Bachelier 的随机游走同时否定了两者。
\textbf{其二，分岔点在时间尺度}。1985 年前后同时出现的过度反应证据与动量证据，其差异并不在方法或数据，而在形成期与持有期的长度。
\textbf{其三，批判的降临是同步的}。数据窥探、发表后衰减、拥挤与尾部风险，对两者构成同构的威胁。

这三个事实合起来，支撑了引言中的论断：讲“两策略对立”，是把同一现象的两个切片读成了两种现象。下一节把这个直觉写成数学。

\subsection{一条常被忽略的支线：波动率的均值回归}

在价格方向的争论之外，还有一条几乎平行的支线，它同样古老，却很少被放进“趋势 vs 反转”的比较里。

Engle(1982) 提出的 ARCH 族与 Bollerslev(1986) 的 GARCH 把“波动率自身会回归其长期水平”变成可估计的命题\cite{engle1982,bollerslev1986}。这条支线的经验规律同样清晰：波动率具有显著的自相关（聚类），且在足够长的尺度上回归均值——这正是“波动率均值回归”策略族的基础。此后 Carr 与 Wu(2009) 用指数期权定价出方差风险溢价\cite{carrwu2009}，Bollerslev 等(2009) 发现该溢价能显著预测短期股票收益\cite{bollerslev2009}。

它为什么重要？因为它揭示了一件事：\emph{“均值回归”这个标签下其实装着不止一种载体}。价格的均值回归（$\theta>0$ 的 OU）与波动率的均值回归（GARCH 的条件方差收敛）在数学上不同构——前者的状态变量是一阶矩，后者是二阶矩。这就带来一个操作上的分岔：用价格反转的逻辑去做波动率交易会得到错误结论（波动率不“偏离均值后必然回归”到同一速度），反之亦然。把两者混淆，是实务中一类常见错误的来源。

这条支线也提供了一个对照视角：价格方向上的“趋势 vs 反转”之争，在波动率维度上并不存在对称的争论——因为波动率均值回归的机制（风险溢价的期限结构）比价格方向的机制更容易被数据锁定。

\subsection{被后来文献反复重写的三个节点}

把这条编年史压缩，有三个节点经历了最多的重写，它们分别对应本文后面三章的核心争论。

\textbf{技术规则是否有效（1966--1999）}：Fama 与 Blume(1966) 判定过滤规则无效\cite{famablume1966}，Brock 等(1992) 推翻这一判定\cite{brock1992}，Sullivan 等(1999) 再用数据窥探校正把结论拉回“样本外不显著”\cite{sullivan1999}。同一问题被改了三次，留下的教训是：\emph{“技术规则有效吗”这个问法本身不完整，必须附带“在多大的搜索空间内、用哪种多重检验校正”}。

\textbf{长周期反转是否真实（1985--1993）}：De Bondt--Thaler(1985)、Fama--French(1988)、Poterba--Summers(1988) 给出肯定\cite{debondt1985,famafrench1988,poterbsummers1988}，Richardson(1993) 系统指出这些证据对小样本与检验设定高度敏感\cite{richardson1993}。争论至今没有闭合，但它把“检验功效”概念引入了金融时间序列研究。

\textbf{时间序列动量是否稳健（2012--2020）}：Moskowitz 等(2012) 的 58 合约证据\cite{mop2012}与 Hurst 等(2017) 的百年样本\cite{hop2017}，对上 Huang 等(2020) 的逐资产再检验\cite{huang2020}与 Pitkäjärvi 等(2020) 的跨资产质疑\cite{pitkajarvi2020}。分歧的实质见第~\ref{sec:confidence}~节：它不是“有效或无效”的二分，而是“收益来自预测还是来自缩放”的分解问题。

三个节点的共同形态是同一件事：\emph{一个看似干净的实证结论，在被追问“搜索空间多大、样本多长、成本多少”之后，会逐步失去它的确定性}。这正是本文设立可信度分级的直接理由。
